\documentclass[11pt]{article}
\usepackage[margin=1in]{geometry}
\usepackage{amsmath,amssymb,amsthm}
\usepackage{hyperref}
\usepackage{enumitem}

\newtheorem{theorem}{Theorem}
\newtheorem{lemma}[theorem]{Lemma}
\newtheorem{corollary}[theorem]{Corollary}
\theoremstyle{definition}
\newtheorem{definition}[theorem]{Definition}
\theoremstyle{remark}
\newtheorem{remark}[theorem]{Remark}

\usepackage{xurl}
\let\oldus\_
\renewcommand{\_}{\oldus\allowbreak}
\newcommand{\e}{\operatorname{e}}

\title{A disproof of Conjecture 00000001354\\[2pt]
\large Square-root-phase prime sums admit no uniform bound $x^{1/4+\varepsilon}$}
\date{}

\begin{document}
\sloppy
\maketitle

\section{The conjecture}

Conjecture 00000001354 reads (English version):

\begin{quote}
\emph{Definition:} The square-phase prime sum is $\sum_{p\le x} e(\sqrt p)$. \emph{Conjecture:}
$\sum_{p\le x} e(\sqrt p) \ll x^{1/4+\varepsilon}$; and the same upper bound holds uniformly after
replacing $\sqrt p$ by $\alpha\sqrt p$ for any fixed $0<\alpha\le 1$ -- square-root cancellation for
prime sums with square-root phases, uniformly.
\end{quote}

The Chinese version (in \texttt{conjectures/00000001354.md}) says, in transliteration, ``ba $\sqrt p$
huan wei $\alpha\sqrt p$ (renyi guding $0<\alpha\le1$) hou tong yi shangjie \emph{yizhi chengli}, ji
pingfanggen xiangwei de sushu he you \emph{junyun} pingfanggen dixiao'': ``after replacing $\sqrt p$ by
$\alpha\sqrt p$ (any fixed $0<\alpha\le 1$) the same upper bound \emph{holds uniformly}, i.e.\ prime sums
with square-root phases have \emph{uniform} square-root cancellation''.

Write $e(t) = \exp(2\pi i t)$ and $S_\alpha(x) = \sum_{p\le x} e(\alpha\sqrt p)$, the sum over primes
$p \le x$. The conjecture is the conjunction of
\begin{enumerate}[label=(\arabic*)]
\item for every $\varepsilon>0$: $S_1(x) \ll_\varepsilon x^{1/4+\varepsilon}$;
\item for every $\varepsilon>0$: $S_\alpha(x) \ll_\varepsilon x^{1/4+\varepsilon}$ uniformly in
  $0<\alpha\le1$.
\end{enumerate}

\textbf{Result.} Clause (2) is false for every $0<\varepsilon<3/4$, so the conjecture is false. Taking
$\alpha = 1/(8\sqrt x)$ makes every phase lie in $[0,\pi/4]$, so $|S_\alpha(x)| \ge \cos(\pi/4)\,\pi(x)$,
and Chebyshev's lower bound for $\pi(x)$ beats $C x^{1/4+\varepsilon}$ for every $C$. Everything is
formalized in Lean~4 with Mathlib (Section~\ref{sec:lean}).

\section{Reading}\label{sec:reading}

\begin{enumerate}[label=(\alph*)]
\item \textbf{Uniformity.} A bound $f(x,\alpha) \ll g(x)$ holds \emph{uniformly} in $\alpha \in A$ if
  one implied constant (and one range of $x$) serves every $\alpha \in A$; when the constant may depend
  on a parameter, this is indicated by a subscript (Section~\ref{sec:sources}). We formalize clause (2)
  as: for every $\varepsilon > 0$ there are $C$ and $X$, independent of $\alpha$, with
  $|S_\alpha(x)| \le C x^{1/4+\varepsilon}$ for all $\alpha\in(0,1]$ and all $x \ge X$. Allowing a
  threshold $X$ only weakens the claim compared with the usual convention ($x\ge 2$), so the refutation
  covers both.
\item \textbf{Why this is the literal reading.} The Chinese text says the bound \emph{yizhi chengli}
  (holds uniformly), and then restates the claim as \emph{junyun} (uniform) square-root cancellation.
  \emph{Yizhi} is the standard Chinese term for ``uniform(ly)'' in analysis (\emph{yizhi shoulian},
  uniform convergence; \emph{yizhi lianxu}, uniform continuity). The English says ``uniformly'' twice.
  If the constant could depend on $\alpha$, both words would carry no content, because ``the same upper
  bound'' already fixes the exponent $1/4+\varepsilon$. The phrase ``for any fixed $0<\alpha\le1$''
  names the parameter range over which uniformity is asserted, as in the standard phrasing ``for fixed
  $\alpha$ \dots, uniformly in $\alpha$''.
\item \textbf{What is not claimed.} We do not refute clause (1) on its own, nor the non-uniform
  reading in which $C$ may depend on $\alpha$. Both are open problems about primes in short intervals,
  and nothing below depends on them. The conjecture, read with its uniform clause, is false because
  that clause is false.
\end{enumerate}

\section{Definitions and sources}\label{sec:sources}

\begin{definition}
$e(t) = \exp(2\pi i t)$ for real $t$; $\pi(x)$ is the number of primes $p\le x$; for real $\alpha, x$,
$S_\alpha(x) = \sum_{p\le x,\ p \text{ prime}} e(\alpha\sqrt p)$. The conjecture's square-phase prime
sum is $S_1(x)$.
\end{definition}

Sources (quotations verbatim):
\begin{itemize}
\item Wikipedia, \emph{Big O notation}, \url{https://en.wikipedia.org/wiki/Big_O_notation}, on bounds
  in several variables: ``In all of these examples, the bound is uniform in both variables.'' On a
  parameter-dependent constant, marked by a subscript: ``Here the implied constant depends on the size
  of the domain.''
\item Chebyshev's lower bound, as formalized in Mathlib (file
  \texttt{Mathlib/\allowbreak NumberTheory/\allowbreak Chebyshev.lean},
  theorem \texttt{Chebyshev.pi\_ge'}): for real $x>1$,
  \[ \pi(x) \;\ge\; \frac{(x-1)\log 2 - \log(x+2)}{\log x}. \]
\end{itemize}

\section{The proof}

\begin{lemma}\label{lem:phase}
For $x>0$ and $\alpha = 1/(8\sqrt x)$: $\ \operatorname{Re} S_\alpha(x) \ge \frac{\sqrt2}{2}\,\pi(x)$.
\end{lemma}

\begin{proof}
$\operatorname{Re} e(t) = \cos(2\pi t)$. For a prime $p\le x$, the phase is
$2\pi\alpha\sqrt p = \frac{\pi}{4}\cdot\frac{\sqrt p}{\sqrt x} \in [0,\pi/4]$. Since $\cos$ decreases on
$[0,\pi]$, $\cos(2\pi\alpha\sqrt p) \ge \cos(\pi/4) = \sqrt2/2$. Summing over the $\pi(x)$ primes $p\le x$
gives the claim.
\end{proof}

\begin{lemma}\label{lem:cheb}
Let $0<\theta<1$. For all real $K$ and $X$ there is $x\ge\max(X,2)$ with $K x^\theta < \pi(x)$.
\end{lemma}

\begin{proof}
Let $\eta = (1-\theta)/2 \in (0, 1/2)$ and choose $x \ge \max(X,2)$ with $y := x^\eta > 4(|K|+2)/\eta$
(possible since $x^\eta\to\infty$). Put $z = x^{1-\eta}$, so $x = zy$, $x^\theta y = z$, and $y \le z$
(as $\eta \le 1-\eta$ and $x\ge1$). We use $\log x \le x^\eta/\eta = y/\eta$, $\log(x+2) \le 2\log x$
(since $x+2\le x^2$), $\log 2 > 1/2$ and hence $(x-1)\log 2 \ge x/4$. Chebyshev's bound gives
$(x-1)\log 2 - \log(x+2) \le \pi(x)\log x$. Multiplying by $\eta$ and using the inequalities above,
\[
  \eta x/4 - 2y \;\le\; \eta\bigl((x-1)\log 2 - \log(x+2)\bigr) \;\le\; \eta\,\pi(x)\log x \;\le\; \pi(x)\, y .
\]
Since $\eta y/4 > |K|+2$, we get $\eta x/4 = (\eta y/4)\, z > (|K|+2) z \ge K z + 2y$, and so
$K x^\theta y = K z < \eta x/4 - 2y \le \pi(x) y$. Dividing by $y>0$ gives $K x^\theta < \pi(x)$.
\end{proof}

\begin{theorem}\label{thm:main}
For every $0<\varepsilon<3/4$ and all $C, X$ there are $\alpha\in(0,1]$ and $x\ge X$ with
$|S_\alpha(x)| > C x^{1/4+\varepsilon}$. Hence clause (2) is false, and so is the conjecture.
\end{theorem}

\begin{proof}
Apply Lemma~\ref{lem:cheb} with $\theta = 1/4+\varepsilon$ and $K = 2C$ to get $x \ge \max(X,2)$ with
$2Cx^\theta < \pi(x)$. Put $\alpha = 1/(8\sqrt x)$. Then $0<\alpha\le 1$ because $\sqrt x\ge1$, and by
Lemma~\ref{lem:phase}
\[
 |S_\alpha(x)| \;\ge\; \operatorname{Re}S_\alpha(x) \;\ge\; \tfrac{\sqrt2}{2}\,\pi(x)
 \;\ge\; \tfrac12\,\pi(x) \;>\; C x^{1/4+\varepsilon}.
\]
(For $\varepsilon \ge 3/4$ the bound is trivial, since $|S_\alpha(x)| \le \pi(x) \le x$.)
\end{proof}

\section{Lean 4 formalization}\label{sec:lean}

The directory \texttt{lean/} contains a Lean~4 project (Lean \texttt{v4.33.1}, Mathlib \texttt{v4.33.1}).
\texttt{Conjecture1354.lean} imports the single source file \texttt{Conjecture1354/Basic.lean}
(179 lines, namespace \texttt{Conjecture1354}, only import \texttt{Mathlib}).

\paragraph{Definitions.}
\begin{itemize}
\item \texttt{e t := Complex.exp (2 * $\pi$ * Complex.I * t)}.
\item \texttt{S $\alpha$ x := $\sum$ p $\in$ Nat.primesLE $\lfloor$x$\rfloor_+$, e ($\alpha$ * $\sqrt{}$(p : $\mathbb{R}$))};
  Mathlib's \texttt{Nat.primesLE n} is the finset of primes $\le n$, so this is $S_\alpha(x)$.
\item \texttt{AlphaOneClause}: \texttt{$\forall$ $\varepsilon$ > 0, $\exists$ C X, $\forall$ x $\ge$ X,
  $\|$S 1 x$\|$ $\le$ C * x \^{} (1/4 + $\varepsilon$)} (clause (1)).
\item \texttt{UniformBound $\varepsilon$}: \texttt{$\exists$ C X, $\forall$ $\alpha$ $\in$ Set.Ioc 0 1,
  $\forall$ x $\ge$ X, $\|$S $\alpha$ x$\|$ $\le$ C * x \^{} (1/4 + $\varepsilon$)}.
\item \texttt{UniformClause := $\forall$ $\varepsilon$ > 0, UniformBound $\varepsilon$} (clause (2)).
\item \texttt{Conjecture := AlphaOneClause $\wedge$ UniformClause}.
\end{itemize}

\paragraph{Theorems.}
\begin{itemize}
\item \texttt{e\_re : (e t).re = Real.cos (2 * $\pi$ * t)}.
\item \texttt{re\_S\_ge (hx : 0 < x) : $\sqrt{2}$ / 2 * $\pi$($\lfloor$x$\rfloor_+$) $\le$ (S (1 / (8 * $\sqrt{}$x)) x).re}
  (Lemma~\ref{lem:phase}; $\pi$ is \texttt{Nat.primeCounting}).
\item \texttt{exists\_large (h$\theta$0 : 0 < $\theta$) (h$\theta$1 : $\theta$ < 1) (K X : $\mathbb{R}$) :
  $\exists$ x, X $\le$ x $\wedge$ 2 $\le$ x $\wedge$ K * x \^{} $\theta$ < $\pi$($\lfloor$x$\rfloor_+$)}
  (Lemma~\ref{lem:cheb}, from \texttt{Chebyshev.pi\_ge'} and \texttt{Real.log\_le\_rpow\_div}).
\item \texttt{not\_uniformBound (h$\varepsilon$0 : 0 < $\varepsilon$) (h$\varepsilon$ : $\varepsilon$ < 3 / 4) :
  $\neg$ UniformBound $\varepsilon$} (Theorem~\ref{thm:main}).
\item \texttt{not\_uniformClause : $\neg$ UniformClause}.
\item The main theorem
\[
  \texttt{conjecture1354\_false} :\ \neg\,\texttt{Conjecture}.
\]
\end{itemize}

\paragraph{Axiom audit.} There is no \texttt{sorry}, no \texttt{native\_decide} and no added axiom.
\texttt{\#print axioms} reports
\begin{verbatim}
'Conjecture1354.conjecture1354_false' depends on axioms:
  [propext, Classical.choice, Quot.sound]
\end{verbatim}
and the same for \texttt{not\_uniformBound}. The file also elaborates with
\texttt{-DwarningAsError=true}. To check:
\begin{verbatim}
cd lean
lake exe cache get
lake build
\end{verbatim}

\end{document}
